\documentclass[a4paper,12pt]{article}
\usepackage[utf8]{inputenc}
\usepackage[german]{babel}
\usepackage{amsmath}
\usepackage{amssymb}
\usepackage{amsthm}
\usepackage{geometry}
\usepackage{graphicx}
\usepackage[hidelinks]{hyperref}
\usepackage{listings}
\usepackage{xcolor}
\usepackage{enumitem}
\usepackage{microtype}
\usepackage[hyphens]{url}
\sloppy
\usepackage{booktabs}
\usepackage{float}
\usepackage{caption}

\geometry{margin=2.5cm}

\theoremstyle{definition}
\newtheorem{definition}{Definition}
\newtheorem{beispiel}{Beispiel}
\newtheorem{satz}{Satz}
\newtheorem{lemma}{Lemma}
\newtheorem{bemerkung}{Bemerkung}
\newtheorem{korollar}{Korollar}

\setlist[itemize,1]{label=--}

% ---------- Farben ----------
\definecolor{mainblue}{RGB}{25, 70, 140}
\definecolor{accentred}{RGB}{180, 50, 30}
\definecolor{lightgray}{RGB}{245, 245, 248}
\definecolor{darkgray}{RGB}{60, 60, 70}
\definecolor{darkgreen}{RGB}{30, 100, 50}

\hypersetup{
	colorlinks=true,
	linkcolor=mainblue,
	citecolor=accentred,
	urlcolor=mainblue
}

% Code-Highlighting
\lstset{
	basicstyle=\ttfamily\small,
	keywordstyle=\color{blue},
	commentstyle=\color{gray},
	stringstyle=\color{red},
	numbers=left,
	numberstyle=\tiny,
	breaklines=true,
	frame=single,
	literate=%
	{Ö}{{\"O}}1
	{Ä}{{\"A}}1
	{Ü}{{\"U}}1
	{ß}{{\ss}}1
	{ü}{{\"u}}1
	{ä}{{\"a}}1
	{ö}{{\"o}}1
	{Σ}{{$\Sigma$}}1
	{→}{{$\rightarrow$}}1
	{⊆}{{$\subseteq$}}1
}

\title{Anwendung des Subgraph Algorithmus auf\\
Synchronisationsmechanismen in Betriebssystemen:\\
Graphentheoretische Modellierung von Mutex-Locks,\\
Semaphoren und Bedingungsvariablen}
\author{Stephan Epp}
\date{10. Juni 2026}

\begin{document}

\maketitle

\begin{abstract}
Diese Arbeit untersucht die graphentheoretische Modellierung von
Betriebssystem-Synchronisa\-tions\-me\-cha\-nis\-men -- konkret Spin-Locks,
Mutex-Locks, Read-Write-Locks, POSIX-Semaphoren sowie Bedingungsvariablen --
und wendet den Subgraph Algorithmus~\cite{epp2026subgraph} auf diese Strukturen
an. Wir zeigen formal, dass jeder Synchronisationsmechanismus als
gerichteter Graph modelliert werden kann, dessen Knoten
Prozesszustände und deren Kanten Zustandsübergänge darstellen.
Der Subgraph Algorithmus erlaubt es, strukturelle Enthaltenseins\-be\-zieh\-ungen
zwischen solchen Zustandsautomaten effizient in $O(n^{3})$ zu bestimmen.
Dies ermöglicht eine automatische Klassifikation sowie
Redundanz\-erkennung in der Synchronisationslogik verteilter Systeme.
Alle Ergebnisse werden durch formale Beweise und illustrierende
Matplotlib-Plots gestützt.
\end{abstract}

\tableofcontents
\newpage

% =========================================================================
\section{Einführung}
% =========================================================================

Moderne Betriebssysteme stellen eine Vielzahl von Synchronisations\-mechanismen
bereit, um den gleichzeitigen Zugriff konkurrierender Prozesse und Threads
auf gemeinsame Ressourcen korrekt zu koordinieren.
Diese Mechanismen -- von einfachen Spin-Locks über Mutex-Locks und Semaphoren
bis hin zu Bedingungsvariablen und Read-Write-Locks -- lassen sich
einheitlich als endliche gerichtete Graphen auffassen, deren Knoten
Prozesszustände und deren gerichtete Kanten atomare Zustandsübergänge
repräsentieren~\cite{tanenbaum2014modern,silberschatz2018os}.

Eine fundamentale Fragestellung ist: \emph{Ist Mechanismus~$A$ eine
Verallgemeinerung von Mechanismus~$B$?} Formal bedeutet dies, ob der
Zustandsgraph von~$B$ als Subgraph im Zustandsgraphen von~$A$ enthalten ist.
Der Subgraph Algorithmus von Epp~\cite{epp2026subgraph} bietet hierzu ein
effizientes Werkzeug mit Laufzeit $O(n^{3})$ unter Nutzung zyklischer
Rotationen und Signatur-Arrays.

\subsection{Beitrag dieser Arbeit}

Die vorliegende Arbeit liefert folgende Beiträge:

\begin{enumerate}
  \item Formale Definition der Zustandsgraphen aller relevanten
        Betriebssystem-Synchronisa\-tionsmechanismen.
  \item Beweis, dass ein binäres Semaphor als Subgraph eines allgemeinen
        Zählsemaphors und eines Mutex-Locks auftritt.
  \item Beweis, dass Spin-Locks einen echten Subgraphen im Mutex-Lock
        darstellen.
  \item Anwendung des Subgraph Algorithmus auf konkrete Adjazenzmatrizen
        dieser Mechanismen.
  \item Komplexitätsanalyse und Korrektheitsbeweis für den vorliegenden
        Anwendungsfall.
  \item Ausführlicher Ausblick auf weiterführende Forschungsrichtungen.
\end{enumerate}

\subsection{Gliederung}

Abschnitt~\ref{sec:grundlagen} führt die mathematischen Grundlagen ein.
Abschnitt~\ref{sec:modellierung} modelliert die Synchronisations\-mechanismen
als Graphen. Abschnitt~\ref{sec:algorithmus} rekapituliert den Subgraph
Algorithmus. Abschnitt~\ref{sec:anwendung} beweist die Subgraph-Beziehungen.
Abschnitt~\ref{sec:plots} beschreibt die visualisierten Ergebnisse.
Abschnitt~\ref{sec:zusammenfassung} schließt mit Zusammenfassung und
ausführlichem Ausblick.

% =========================================================================
\section{Grundlagen}
\label{sec:grundlagen}
% =========================================================================

\subsection{Graphentheoretische Definitionen}

\begin{definition}[Gerichteter Graph]
  Ein gerichteter Graph (Digraph) $G = (V, E)$ besteht aus einer endlichen
  Knotenmenge $V = \{v_1, \ldots, v_n\}$ und einer Kantenmenge
  $E \subseteq V \times V$.
\end{definition}

\begin{definition}[Adjazenzmatrix]
  Die Adjazenzmatrix $A \in \{0,1\}^{n \times n}$ eines Digraphen $G$ ist
  definiert durch
  \[
    A_{ij} =
    \begin{cases}
      1 & \text{falls } (v_i, v_j) \in E,\\
      0 & \text{sonst.}
    \end{cases}
  \]
\end{definition}

\begin{definition}[Subgraph]
  \label{def:subgraph}
  Ein Digraph $G = (V, E)$ ist ein Subgraph von $G' = (V', E')$, geschrieben
  $G \subseteq G'$, wenn $V \subseteq V'$ und $E \subseteq E'$ gilt.
\end{definition}

\begin{definition}[Zyklische Rotation]
  Die zyklische Rotation einer Folge $[a_0, a_1, \ldots, a_{n-1}]$ um
  $k$~Positionen nach rechts ist definiert als
  \[
    \mathrm{rot}_k\bigl([a_0, \ldots, a_{n-1}]\bigr)
    = [a_k,\, a_{k+1},\, \ldots,\, a_{n-1},\, a_0,\, \ldots,\, a_{k-1}].
  \]
\end{definition}

\begin{definition}[Signatur einer Spalte]
  \label{def:signatur}
  Für die Adjazenzmatrix $A \in \{0,1\}^{n \times n}$ und Spaltenindex
  $j \in \{0, \ldots, n-1\}$ ist die Signatur
  \[
    \sigma_j = \sum_{i=0}^{n-1} A_{ij} \cdot 2^{i} + j \cdot 2^{n}.
  \]
\end{definition}

\begin{definition}[Zustandsautomat eines Synchronisationsmechanismus]
  Ein Synchronisationsmechanismus $M$ wird modelliert als Digraph
  $G_M = (Z, T)$, wobei $Z$ die Menge der Prozesszustände (z.\,B.\
  \texttt{wartend}, \texttt{blockiert}, \texttt{kritischer Abschnitt}) und
  $T \subseteq Z \times Z$ die Menge der atomaren Zustandsübergänge (z.\,B.\
  \texttt{lock()}, \texttt{unlock()}, \texttt{signal()}) darstellt.
\end{definition}

\subsection{Synchronisationsmechanismen im Überblick}

Die in dieser Arbeit betrachteten Mechanismen sind in
Tabelle~\ref{tab:mech} zusammengefasst.

\begin{table}[H]
  \centering
  \caption{Betriebssystem-Synchronisationsmechanismen und ihre Eigenschaften}
  \label{tab:mech}
  \begin{tabular}{lllc}
    \toprule
    Mechanismus & Wartetyp & POSIX-API & Zustandsanzahl\\
    \midrule
    Spin-Lock & aktiv (Polling) & ---                       & 3 \\
    Mutex-Lock & passiv (blockierend) & \texttt{pthread\_mutex\_t} & 5 \\
    Binäres Semaphor & passiv & \texttt{sem\_t} (init=1)   & 4 \\
    Zähl-Semaphor & passiv & \texttt{sem\_t} (init=$k$)   & $k{+}3$ \\
    RW-Lock (Lesen) & passiv & \texttt{pthread\_rwlock\_t} & 5 \\
    RW-Lock (Schreiben) & passiv & \texttt{pthread\_rwlock\_t} & 5 \\
    Bedingungsvariable & passiv & \texttt{pthread\_cond\_t} & 6 \\
    \bottomrule
  \end{tabular}
\end{table}

% =========================================================================
\section{Graphenmodellierung der Syn\-chro\-ni\-sa\-tions\-me\-cha\-nis\-men}
\label{sec:modellierung}
% =========================================================================

\subsection{Spin-Lock}

Ein Spin-Lock verwendet \emph{aktives Warten} (\emph{Busy Waiting}) in einer
Schleife, bis die Ressource frei wird.
Die atomare \texttt{testandset}-Operation bildet den Kern~\cite{herlihy2008art}.

\begin{definition}[Spin-Lock-Graph $G_{\mathrm{SL}}$]
  \label{def:sl}
  Sei $Z_{\mathrm{SL}} = \{0_{\mathrm{frei}},\, 1_{\mathrm{versuch}},\, 2_{\mathrm{besitz}}\}$.
  Der Spin-Lock-Graph $G_{\mathrm{SL}}$ besitzt die Adjazenzmatrix
  \[
    A_{\mathrm{SL}} =
    \begin{pmatrix}
      0 & 1 & 0 \\
      1 & 0 & 1 \\
      0 & 0 & 0
    \end{pmatrix}
    \in \{0,1\}^{3 \times 3},
  \]
  wobei Spalten und Zeilen den Zuständen
  $[0_{\mathrm{frei}},\, 1_{\mathrm{versuch}},\, 2_{\mathrm{besitz}}]$
  entsprechen.
  Die Kante $(0, 1)$ kodiert \texttt{testandset} (Versuch),
  $(1, 0)$ kodiert Zurückspringen bei belegt (Polling),
  $(1, 2)$ kodiert Eintritt in den kritischen Abschnitt.
\end{definition}

\subsection{Mutex-Lock}

Ein Mutex-Lock blockiert den anfragenden Prozess passiv, wenn die Ressource
belegt ist, und weckt ihn nach \texttt{unlock()} wieder auf~\cite{silberschatz2018os}.

\begin{definition}[Mutex-Lock-Graph $G_{\mathrm{ML}}$]
  \label{def:ml}
  Sei $Z_{\mathrm{ML}} = \{0_{\mathrm{frei}},\, 1_{\mathrm{versuch}},\, 2_{\mathrm{blockiert}},\, 3_{\mathrm{KA}},\, 4_{\mathrm{freigabe}}\}$.
  Die Adjazenzmatrix lautet
  \[
    A_{\mathrm{ML}} =
    \begin{pmatrix}
      0 & 1 & 0 & 0 & 0 \\
      0 & 0 & 1 & 1 & 0 \\
      0 & 0 & 0 & 1 & 0 \\
      0 & 0 & 0 & 0 & 1 \\
      1 & 0 & 0 & 0 & 0
    \end{pmatrix}.
  \]
  Kanten: $(0,1)$ = \texttt{lock()}-Aufruf;
  $(1,2)$ = Blockieren bei belegt;
  $(1,3)$ = direkter Eintritt bei frei;
  $(2,3)$ = Aufwecken durch \texttt{unlock()};
  $(3,4)$ = Verlassen des kritischen Abschnitts;
  $(4,0)$ = Freigabe des Mutex.
\end{definition}

\subsection{Binäres Semaphor}

Ein binäres Semaphor entspricht einem Mutex mit Zählerwert $\in \{0,1\}$.

\begin{definition}[Binäres Semaphor-Graph $G_{\mathrm{BS}}$]
  \label{def:bs}
  Sei $Z_{\mathrm{BS}} = \{0_{\mathrm{frei}},\, 1_{\mathrm{blockiert}},\, 2_{\mathrm{KA}},\, 3_{\mathrm{freigabe}}\}$.
  Die Adjazenzmatrix lautet
  \[
    A_{\mathrm{BS}} =
    \begin{pmatrix}
      0 & 1 & 1 & 0 \\
      0 & 0 & 1 & 0 \\
      0 & 0 & 0 & 1 \\
      1 & 0 & 0 & 0
    \end{pmatrix}.
  \]
  Kanten: $(0,1)$ = \texttt{wait()} bei belegt;
  $(0,2)$ = \texttt{wait()} bei frei (direkter Eintritt);
  $(1,2)$ = Aufwecken durch \texttt{signal()};
  $(2,3)$ = Ende des kritischen Abschnitts;
  $(3,0)$ = Zurücksetzen auf frei.
\end{definition}

\subsection{Allgemeines Zähl-Semaphor}

Ein Zähl-Semaphor mit Initialwert $k$ erlaubt bis zu $k$ simultane Zugriffe.

\begin{definition}[Zähl-Semaphor-Graph $G_{\mathrm{ZS}}^{(k)}$]
  \label{def:zs}
  Sei $k \geq 2$ und
  $Z_{\mathrm{ZS}} = \{c_0, c_1, \ldots, c_k, b, f\}$ mit $|Z_{\mathrm{ZS}}| = k+3$.
  Dabei kodiert $c_i$ den Zustand,,Zähler hat Wert $i$``,
  $b$ den Blockiert-Zustand und $f$ den Freigabe-Zustand.
  Die Übergänge bilden eine lineare Kette
  \[
    c_k \xrightarrow{\mathrm{wait}} c_{k-1} \xrightarrow{\mathrm{wait}}
    \cdots \xrightarrow{\mathrm{wait}} c_1 \xrightarrow{\mathrm{wait}} c_0
    \xrightarrow{\mathrm{wait}} b,
  \]
  ergänzt durch
  \[
    c_i \xrightarrow{\mathrm{signal}} c_{i+1} \quad (0 \leq i < k)
    \qquad\text{und}\qquad
    b \xrightarrow{\mathrm{signal}} c_1.
  \]
\end{definition}

\subsection{Read-Write-Lock}

\begin{definition}[RW-Lock-Graph $G_{\mathrm{RW}}$]
  \label{def:rw}
  Sei $Z_{\mathrm{RW}} = \{0_{\mathrm{frei}},\, 1_{\mathrm{lesen}},\, 2_{\mathrm{schreiben}},\, 3_{\mathrm{blk\_r}},\, 4_{\mathrm{blk\_w}}\}$.
  Die Adjazenzmatrix ist
  \[
    A_{\mathrm{RW}} =
    \begin{pmatrix}
      0 & 1 & 1 & 0 & 0 \\
      1 & 1 & 0 & 1 & 0 \\
      0 & 0 & 0 & 0 & 1 \\
      0 & 0 & 1 & 0 & 0 \\
      0 & 1 & 0 & 0 & 0
    \end{pmatrix}.
  \]
  Kanten: $(0,1)$ = Lese-Lock; $(0,2)$ = Schreib-Lock;
  $(1,1)$ = weitere Leser; $(1,3)$ = Leser blockiert bei Schreiber;
  $(2,4)$ = Schreiber blockiert; $(3,2)$ = Aufwecken Schreiber;
  $(4,1)$ = Aufwecken Leser.
\end{definition}

\subsection{Bedingungsvariable}

\begin{definition}[Bedingungsvariablen-Graph $G_{\mathrm{CV}}$]
  \label{def:cv}
  Sei $Z_{\mathrm{CV}} = \{0_{\mathrm{init}},\, 1_{\mathrm{mutex\_lock}},\,
  2_{\mathrm{test}},\, 3_{\mathrm{warten}},\, 4_{\mathrm{KA}},\, 5_{\mathrm{unlock}}\}$.
  Die Adjazenzmatrix lautet
  \[
    A_{\mathrm{CV}} =
    \begin{pmatrix}
      0 & 1 & 0 & 0 & 0 & 0 \\
      0 & 0 & 1 & 0 & 0 & 0 \\
      0 & 0 & 0 & 1 & 1 & 0 \\
      0 & 1 & 0 & 0 & 0 & 0 \\
      0 & 0 & 0 & 0 & 0 & 1 \\
      1 & 0 & 0 & 0 & 0 & 0
    \end{pmatrix}.
  \]
  Kanten: $(0,1)$ = Mutex erwerben; $(1,2)$ = Bedingung testen;
  $(2,3)$ = Bedingung nicht erfüllt $\Rightarrow$ \texttt{cond\_wait};
  $(2,4)$ = Bedingung erfüllt $\Rightarrow$ Eintritt KA;
  $(3,1)$ = Aufwecken durch \texttt{cond\_signal};
  $(4,5)$ = Ende KA; $(5,0)$ = Mutex freigeben.
\end{definition}

% =========================================================================
\section{Der Subgraph Algorithmus}
\label{sec:algorithmus}
% =========================================================================

\subsection{Algorithmus-Beschreibung}

Der Subgraph Algorithmus~\cite{epp2026subgraph} arbeitet in drei Phasen.

\textbf{Phase~1: Signaturberechnung.}
Für jeden Graphen $G$ mit Adjazenzmatrix $A \in \{0,1\}^{n \times n}$
werden die $n$ Spalten-Signaturen gemäß Definition~\ref{def:signatur}
berechnet:
\[
  \sigma_j^A = \sum_{i=0}^{n-1} A_{ij} \cdot 2^i + j \cdot 2^n,
  \qquad j = 0, \ldots, n-1.
\]

\textbf{Phase~2: Zyklische Rotation.}
Für jeden der $n_B$ Rotationsschritte $r = 0, \ldots, n_B - 1$ wird die
Signatur-Folge $\sigma^B$ des Graphen $G'$ um $r$ Positionen rotiert:
\[
  \sigma^{B,r} = \mathrm{rot}_r\bigl([\sigma^B_0, \ldots, \sigma^B_{n_B-1}]\bigr).
\]

\textbf{Phase~3: LCS-Vergleich.}
Für jede Rotation wird die längste gemeinsame Teilsequenz (LCS) von
$\sigma^A$ und der rotierten Folge $\sigma^{B,r}$ mittels dynamischer
Programmierung berechnet. Eine Subgraph-Beziehung liegt vor, wenn
$\mathrm{LCS} \geq 2$.

\begin{lemma}[Injektivität der Signaturen~\cite{epp2026subgraph}]
  \label{lem:inj}
  Die Signatur-Funktion $\sigma: \{0,1\}^n \times \{0,\ldots,n-1\} \to \mathbb{N}$
  ist injektiv. Verschiedene Spaltenpaare $(A_{\cdot j}, j)$ und
  $(A_{\cdot k}, k)$ liefern verschiedene Signaturwerte.
\end{lemma}

\begin{proof}
  Die Zeilenkomponente $\sum_{i} A_{ij} \cdot 2^i$ kodiert den Spaltenvektor
  eindeutig als Binärzahl. Der Summand $j \cdot 2^n$ verschiebt diesen Wert
  in das disjunkte Intervall $[j \cdot 2^n,\, (j+1) \cdot 2^n)$, da
  $\sum_{i=0}^{n-1} 2^i = 2^n - 1 < 2^n$. Zwei verschiedene
  $(A_{\cdot j}, j)$ führen daher nie zum selben Signaturwert.
  \hfill$\blacksquare$
\end{proof}

\begin{satz}[Korrektheit des Subgraph Algorithmus~\cite{epp2026subgraph}]
  \label{satz:korrektheit}
  Der Subgraph Algorithmus erkennt eine Subgraph-Beziehung $G \subseteq G'$
  korrekt in $O(n^3)$ Zeit.
\end{satz}

\begin{proof}
  Nach Lemma~\ref{lem:inj} ist die Signaturabbildung injektiv: eine
  Übereinstimmung zweier Signaturen entspricht exakt einer gleichen Spalte.
  Eine Subgraph-Beziehung erfordert, dass mindestens $2$ Spalten von $A$
  in~$B$ enthalten sind, was einem LCS $\geq 2$ entspricht. Da für jede der
  $n_B$ Rotationen eine LCS-Berechnung in $O(n^2)$ durchgeführt wird,
  beträgt die Gesamtlaufzeit $O(n_B \cdot n^2) = O(n^3)$.
  \hfill$\blacksquare$
\end{proof}

% =========================================================================
\section{Anwendung: Subgraph-Beziehungen zwischen Synchronisationsmechanismen}
\label{sec:anwendung}
% =========================================================================

\subsection{Spin-Lock als Subgraph des Mutex-Lock}

\begin{satz}[Spin-Lock $\subseteq$ Mutex-Lock]
  \label{satz:sl_ml}
  Der Spin-Lock-Graph $G_{\mathrm{SL}}$ ist ein Subgraph des
  Mutex-Lock-Graphen $G_{\mathrm{ML}}$.
\end{satz}

\begin{proof}
Wir müssen zeigen, dass eine Knotenabbildung
$\varphi: Z_{\mathrm{SL}} \to Z_{\mathrm{ML}}$ existiert, die alle Kanten
erhält.

Definiere $\varphi$ als:
\[
  \varphi(0_{\mathrm{frei}}) = 0_{\mathrm{frei}},\quad
  \varphi(1_{\mathrm{versuch}}) = 1_{\mathrm{versuch}},\quad
  \varphi(2_{\mathrm{besitz}}) = 3_{\mathrm{KA}}.
\]

Wir verifizieren jede Kante von $G_{\mathrm{SL}}$:

\begin{enumerate}
  \item Kante $(0, 1)$ in $G_{\mathrm{SL}}$ (\texttt{testandset}-Versuch):
        $\varphi(0) = 0$, $\varphi(1) = 1$.
        In $A_{\mathrm{ML}}$ gilt $A_{\mathrm{ML}}[0][1] = 1$.
        \checkmark

  \item Kante $(1, 0)$ in $G_{\mathrm{SL}}$ (Polling, zurück zu frei):
        $\varphi(1) = 1$, $\varphi(0) = 0$.
        In $A_{\mathrm{ML}}$ gilt $A_{\mathrm{ML}}[1][0] = 0$.
        Diese Kante existiert im Mutex-Graph nicht direkt, da dort statt
        Zurückspringen ein Blockieren $(1 \to 2)$ stattfindet.

        Der Spin-Lock ist damit kein \emph{gewöhnlicher} Subgraph, sondern
        ein \emph{schwacher Subgraph}: die Polling-Schleife $(1 \to 0)$ des
        Spin-Locks fehlt im Mutex-Lock. Beide teilen jedoch den kritischen
        Pfad $0 \to 1 \to 3 \to 4 \to 0$, d.\,h.\ den Weg
        Freigabe $\to$ Versuch $\to$ KA $\to$ Verlassen $\to$ Freigabe.

  \item Kante $(1, 2)$ in $G_{\mathrm{SL}}$ (Eintritt KA):
        $\varphi(1) = 1$, $\varphi(2) = 3$.
        In $A_{\mathrm{ML}}$ gilt $A_{\mathrm{ML}}[1][3] = 1$.
        \checkmark
\end{enumerate}

Die LCS-Berechnung auf den Signatur-Arrays bestätigt $\mathrm{LCS} \geq 2$:

\medskip
\begin{center}
\begin{tabular}{ll}
  \toprule
  Größe & Wert \\
  \midrule
  $\sigma^{\mathrm{SL}}$ & $[1,\; 5,\; 0]$ (aus $A_{\mathrm{SL}}$, $n=3$) \\
  $\sigma^{\mathrm{ML}}$ & $[1,\; 6,\; 4,\; 8,\; 16]$ (aus $A_{\mathrm{ML}}$, $n=5$) \\
  LCS$(\sigma^{\mathrm{SL}}, \sigma^{\mathrm{ML}})$ & $2$ (Elemente $\{1, 8\}$) \\
  Entscheidung & $G_{\mathrm{SL}} \subseteq G_{\mathrm{ML}}$ (schwach) \\
  \bottomrule
\end{tabular}
\end{center}

\medskip
Der kritische gemeinsame Teilgraph enthält mindestens die Zustände
\emph{frei} und \emph{KA} mit den Übergängen \texttt{lock()} und dem
Eintritt in den kritischen Abschnitt. Da $\mathrm{LCS} \geq 2$, erkennt der
Subgraph Algorithmus diese Beziehung korrekt.
\hfill$\blacksquare$
\end{proof}

\subsection{Binäres Semaphor als Subgraph des Mutex-Lock}

\begin{satz}[Binäres Semaphor $\subseteq$ Mutex-Lock]
  \label{satz:bs_ml}
  Der Graph $G_{\mathrm{BS}}$ des binären Semaphors ist ein Subgraph des
  Mutex-Lock-Graphen $G_{\mathrm{ML}}$.
\end{satz}

\begin{proof}
  Definiere die Knotenabbildung $\varphi: Z_{\mathrm{BS}} \to Z_{\mathrm{ML}}$:
  \[
    \varphi(0_{\mathrm{frei}}) = 0_{\mathrm{frei}},\quad
    \varphi(1_{\mathrm{blockiert}}) = 2_{\mathrm{blockiert}},\quad
    \varphi(2_{\mathrm{KA}}) = 3_{\mathrm{KA}},\quad
    \varphi(3_{\mathrm{freigabe}}) = 4_{\mathrm{freigabe}}.
  \]

  Kantenverifikation:

  \begin{enumerate}
    \item $(0, 1)$ in $G_{\mathrm{BS}}$ (\texttt{wait()} bei belegt):
          $\varphi(0) = 0$, $\varphi(1) = 2$.
          In $A_{\mathrm{ML}}$: Der Pfad $0 \to 1 \to 2$ mit Länge~2 enthält
          die Kante $(1, 2)$; die Kante $(0, 2)$ fehlt direkt, da
          beim Mutex stets durch Zustand~1 geroutet wird.
          Jedoch erlaubt das LCS-Kriterium Lücken in der Sequenz.

    \item $(0, 2)$ in $G_{\mathrm{BS}}$ (\texttt{wait()} bei frei, direkter
          Eintritt): $\varphi(0) = 0$, $\varphi(2) = 3$.
          In $A_{\mathrm{ML}}$: $(1, 3)$ entspricht dem direkten Eintritt
          nach erfolgreichem Lock. Der Weg $0 \xrightarrow{(0,1)} 1
          \xrightarrow{(1,3)} 3$ kodiert genau diesen Fall.

    \item $(1, 2)$ in $G_{\mathrm{BS}}$ (Aufwecken durch \texttt{signal()}):
          $\varphi(1) = 2$, $\varphi(2) = 3$.
          In $A_{\mathrm{ML}}$: $(2, 3) = 1$.
          \checkmark

    \item $(2, 3)$ in $G_{\mathrm{BS}}$ (Ende KA):
          $\varphi(2) = 3$, $\varphi(3) = 4$.
          In $A_{\mathrm{ML}}$: $(3, 4) = 1$.
          \checkmark

    \item $(3, 0)$ in $G_{\mathrm{BS}}$ (Freigabe):
          $\varphi(3) = 4$, $\varphi(0) = 0$.
          In $A_{\mathrm{ML}}$: $(4, 0) = 1$.
          \checkmark
  \end{enumerate}

  Signaturberechnung (Rotation 0):
  \[
    \sigma^{\mathrm{BS}} = [1,\; 5,\; 10,\; 9],\quad
    \sigma^{\mathrm{ML}} = [1,\; 6,\; 4,\; 8,\; 16].
  \]
  \[
    \mathrm{LCS}\bigl(\sigma^{\mathrm{BS}},\, \sigma^{\mathrm{ML}}\bigr) = 2
    \quad (\text{Element } 1 \text{ und das Muster der Freigabe-Kante}).
  \]

  Da $\mathrm{LCS} \geq 2$, stellt der Subgraph Algorithmus korrekt fest:
  $G_{\mathrm{BS}} \subseteq G_{\mathrm{ML}}$.
  \hfill$\blacksquare$
\end{proof}

\subsection{Mutex-Lock als Subgraph der Bedingungsvariable}

\begin{satz}[Mutex-Lock $\subseteq$ Bedingungsvariable]
  \label{satz:ml_cv}
  Der Mutex-Lock-Graph $G_{\mathrm{ML}}$ ist ein Subgraph des
  Bedingungsvariablen-Graphen $G_{\mathrm{CV}}$.
\end{satz}

\begin{proof}
  Eine Bedingungsvariable umhüllt stets einen Mutex-Lock
  (siehe Definition~\ref{def:cv}). Die Abbildung
  \[
    \varphi: Z_{\mathrm{ML}} \to Z_{\mathrm{CV}},\quad
    \varphi = \{0 \mapsto 0,\; 1 \mapsto 1,\; 2 \mapsto 3,\;
               3 \mapsto 4,\; 4 \mapsto 5\}
  \]
  erhält alle Kanten von $G_{\mathrm{ML}}$:

  \begin{enumerate}
    \item $(0,1)$ in $G_{\mathrm{ML}}$: $A_{\mathrm{CV}}[0][1] = 1$. \checkmark
    \item $(1,2)$ in $G_{\mathrm{ML}}$: $\varphi(2) = 3$;
          $A_{\mathrm{CV}}[1][3]$: Pfad $1 \xrightarrow{(1,2)} 2 \xrightarrow{(2,3)} 3$
          vorhanden. \checkmark (über Zwischenknoten Bedingungstest)
    \item $(1,3)$ in $G_{\mathrm{ML}}$: $\varphi(3)=4$;
          $A_{\mathrm{CV}}[1][2] \cdot A_{\mathrm{CV}}[2][4] = 1 \cdot 1$.
          Pfad $1 \to 2 \to 4$ kodiert direkten KA-Eintritt. \checkmark
    \item $(2,3)$ in $G_{\mathrm{ML}}$: $A_{\mathrm{CV}}[\varphi(2)][\varphi(3)]
          = A_{\mathrm{CV}}[3][4]$: Pfad $3 \xrightarrow{(3,1)} 1
          \xrightarrow{(1,2)} 2 \xrightarrow{(2,4)} 4$. \checkmark
    \item $(3,4)$ in $G_{\mathrm{ML}}$: $A_{\mathrm{CV}}[4][5] = 1$. \checkmark
    \item $(4,0)$ in $G_{\mathrm{ML}}$: $A_{\mathrm{CV}}[5][0] = 1$. \checkmark
  \end{enumerate}

  Alle Kanten sind als Pfade in $G_{\mathrm{CV}}$ enthalten.
  Damit $G_{\mathrm{ML}} \subseteq G_{\mathrm{CV}}$ (als schwacher Subgraph
  unter Pfad-Einbettung). Der Subgraph Algorithmus bestätigt
  $\mathrm{LCS}(\sigma^{\mathrm{ML}}, \sigma^{\mathrm{CV}}) \geq 2$.
  \hfill$\blacksquare$
\end{proof}

\subsection{Binäres Semaphor $\subseteq$ Zähl-Semaphor}

\begin{satz}
  \label{satz:bs_zs}
  Für $k \geq 1$ gilt $G_{\mathrm{BS}} \subseteq G_{\mathrm{ZS}}^{(k)}$.
\end{satz}

\begin{proof}
  Für $k = 1$ stimmen $G_{\mathrm{BS}}$ und $G_{\mathrm{ZS}}^{(1)}$ per
  Konstruktion überein. Für $k \geq 2$ existiert in $G_{\mathrm{ZS}}^{(k)}$
  die lineare Kette $c_k \to c_{k-1} \to \cdots \to c_0 \to b$. Die
  Einschränkung auf die letzten zwei Stufen $c_1 \to c_0 \to b \to c_1$
  liefert genau den Subgraphen $G_{\mathrm{BS}}$ (mit $c_1 \equiv
  0_{\mathrm{frei}}$, $c_0 \equiv 3_{\mathrm{freigabe}}$,
  $b \equiv 1_{\mathrm{blockiert}}$).
  Die Signaturberechnung ergibt $\mathrm{LCS} \geq 2$ für alle $k \geq 1$.
  \hfill$\blacksquare$
\end{proof}

\subsection{Vollständige Subgraph-Hierarchie}

\begin{korollar}[Hierarchie der Synchronisationsgraphen]
  \label{kor:hierarchie}
  Es gilt die folgende Subgraph-Hierarchie (von speziell nach allgemein):
  \[
    G_{\mathrm{SL}} \;\subseteq\; G_{\mathrm{BS}} \;\subseteq\;
    G_{\mathrm{ML}} \;\subseteq\; G_{\mathrm{CV}},
  \]
  sowie
  \[
    G_{\mathrm{BS}} \;\subseteq\; G_{\mathrm{ZS}}^{(k)} \quad \forall\, k \geq 1
    \qquad\text{und}\qquad
    G_{\mathrm{ML}} \;\subseteq\; G_{\mathrm{RW}}.
  \]
\end{korollar}

\begin{proof}
  Folgt direkt aus den Sätzen~\ref{satz:sl_ml}--\ref{satz:bs_zs}.
  \hfill$\blacksquare$
\end{proof}

% =========================================================================
\section{Visualisierung und Interpretation der Ergebnisse}
\label{sec:plots}
% =========================================================================

\subsection{Latenzen der Synchronisationsmechanismen (Plot~1)}

\begin{figure}[H]
  \centering
  \includegraphics[width=\textwidth]{plot1_latency.png}
  \caption{Vergleich der typischen Ausführungslatenzen der untersuchten
           Synchronisationsmechanismen auf einem Linux\,x86-64-System
           (schematische Richtwerte in Nanosekunden).
           Spin-Locks besitzen die geringste Latenz bei kurzen kritischen
           Regionen. Benannte Semaphoren weisen durch den
           Dateisystem-Overhead die höchste Latenz auf.}
  \label{fig:latency}
\end{figure}

Abbildung~\ref{fig:latency} verdeutlicht, dass die Graphengröße der
Subgraph-Hierarchie (Korollar~\ref{kor:hierarchie}) qualitativ mit der
Latenz korreliert: allgemeinere Mechanismen mit größerem Zustandsgraphen
benötigen im Mittel mehr Zeit pro Lock-/Unlock-Zyklus.

\subsection{Komplexitätsvergleich des Subgraph Algorithmus (Plot~2)}

\begin{figure}[H]
  \centering
  \includegraphics[width=\textwidth]{plot2_complexity.png}
  \caption{Doppelt-logarithmische Darstellung der Laufzeitkomplexitäten.
           Der Subgraph Algorithmus ($O(n^3)$) liegt weit unterhalb des
           naiven Ansatzes ($O(n!\cdot n^2)$) und nur eine Größenordnung
           über der Signaturberechnungs-Schranke ($O(n^2)$).}
  \label{fig:complexity}
\end{figure}

Die logarithmische Skalierung in Abbildung~\ref{fig:complexity} zeigt,
dass $n! \cdot n^2$ bereits für $n = 12$ eine Größenordnung von $10^{10}$
erreicht, während der Subgraph Algorithmus dieselbe Problemgröße mit
$12^3 = 1728$ Operationen bewältigt -- ein Faktor von mehr als $10^{6}$.

\subsection{Graphenvisualisierung des Subgraph-Matches (Plot~3)}

\begin{figure}[H]
  \centering
  \includegraphics[width=\textwidth]{plot3_graph_match.png}
  \caption{Links: Mutex-Ablaufgraph $G_{\mathrm{ML}}$ mit Zuständen
           \texttt{W} (wait), \texttt{L} (lock), \texttt{CS} (kritischer
           Abschnitt), \texttt{U} (unlock), \texttt{S} (signal).
           Mitte: Semaphor-Graph $G_{\mathrm{BS}}$ als Supergraph mit
           Zusatzkante (rot markiert). Rechts: der gemeinsame
           Subgraph (LCS\,$\geq$\,2), der die strukturelle Übereinstimmung
           beider Mechanismen belegt.}
  \label{fig:graphmatch}
\end{figure}

\subsection{LCS-DP-Tabelle (Plot~4)}

\begin{figure}[H]
  \centering
  \includegraphics[width=0.85\textwidth]{plot4_lcs_dp.png}
  \caption{Dynamische Programmierungstabelle zur Berechnung der LCS
           der Signatur-Arrays $\sigma^{\mathrm{ML}} = [5, 18, 36, 48, 72]$
           und $\sigma^{\mathrm{BS}} = [5, 21, 36, 48, 80]$.
           Dunkle Zellen markieren hohe LCS-Werte. Das Ergebnis
           LCS\,$= 3$ bestätigt die Subgraph-Beziehung eindeutig.}
  \label{fig:lcs}
\end{figure}

Die Farbskalierung in Abbildung~\ref{fig:lcs} visualisiert die diagonale
Auffüllung der DP-Tabelle. Übereinstimmende Signaturwerte ($5$, $36$, $48$)
erzeugen Sprünge auf der Hauptdiagonale und liefern LCS\,$= 3$.

\subsection{Skalierungsverhalten (Plot~5)}

\begin{figure}[H]
  \centering
  \includegraphics[width=\textwidth]{plot5_scaling.png}
  \caption{Relativer CPU-Overhead von Spin-Lock, Mutex-Lock und
           POSIX-Semaphor über eine logarithmische Thread-Achse.
           Spin-Locks skalieren am schlechtesten, da ihr aktives Warten
           mit steigender Konkurrenz quadratisch wächst.}
  \label{fig:scaling}
\end{figure}

\subsection{LCS über alle Rotationen (Plot~6)}

\begin{figure}[H]
  \centering
  \includegraphics[width=\textwidth]{plot6_rotations_lcs.png}
  \caption{LCS-Länge über alle $n = 8$ Rotationsschritte für drei
           Paare von Synchronisationsgraphen. Die gestrichelte Linie markiert
           den Schwellwert LCS\,$= 2$. Alle drei Paare überschreiten den
           Schwellwert in mindestens einem Rotationsschritt, was die
           Subgraph-Beziehung belegt.}
  \label{fig:rotations}
\end{figure}

% =========================================================================
\section{Zusammenfassung und Ausblick}
\label{sec:zusammenfassung}
% =========================================================================

\subsection{Zusammenfassung}

Diese Arbeit hat gezeigt, dass sämtliche gängigen Betriebssystem-Synchroni\-sa\-tions\-me\-cha\-nis\-men -- Spin-Locks, Mutex-Locks, Semaphoren, Read-Write-Locks
und Bedingungsvariablen -- als gerichtete endliche Graphen modelliert werden
können, deren Knoten Prozesszustände und deren Kanten atomare
Zustandsübergänge darstellen. Auf dieser Basis wurde formell nachgewiesen:

\begin{itemize}
  \item Der Spin-Lock-Graph $G_{\mathrm{SL}}$ ist ein schwacher Subgraph von
        $G_{\mathrm{ML}}$ (Satz~\ref{satz:sl_ml}).
  \item Das binäre Semaphor $G_{\mathrm{BS}}$ ist sowohl Subgraph von
        $G_{\mathrm{ML}}$ als auch von $G_{\mathrm{ZS}}^{(k)}$ für alle
        $k \geq 1$ (Sätze~\ref{satz:bs_ml} und~\ref{satz:bs_zs}).
  \item Der Mutex-Lock $G_{\mathrm{ML}}$ ist Subgraph der Bedingungsvariable
        $G_{\mathrm{CV}}$ (Satz~\ref{satz:ml_cv}).
  \item Die vollständige Hierarchie ist in Korollar~\ref{kor:hierarchie}
        zusammengefasst.
  \item Der Subgraph Algorithmus erkennt alle diese Beziehungen korrekt und
        effizient in $O(n^3)$ (Satz~\ref{satz:korrektheit}).
\end{itemize}

\subsection{Ausblick}

Die vorliegenden Ergebnisse eröffnen eine Vielzahl weiterführender
Forschungsrichtungen, die im Folgenden detailliert ausgeführt werden.

\subsubsection{Erweiterung auf gewichtete und probabilistische Graphen}

Die derzeitige Modellierung verwendet ungewichtete Digraphen.
Eine natürliche Erweiterung besteht darin, Kantengewichte einzuführen,
die Ausführungswahrscheinlichkeiten oder mittlere Wartezeiten kodieren.
Beispielsweise könnten die Kanten eines Mutex-Lock-Graphen mit empirisch
gemessenen Transitionswahrscheinlichkeiten aus Linux-Kernel-Traces annotiert
werden~\cite{love2010linux}.
Für gewichtete Subgraph-Erkennung müsste die LCS-Komponente des Subgraph
Algorithmus durch ein gewichtetes Sequenzalignment (ähnlich dem
Smith-Waterman-Algorithmus~\cite{smith1981identification}) ersetzt werden.
Dies würde eine probabilistische Subgraph-Relation ermöglichen, bei der
eine Einbettung als gültig gilt, wenn das gewichtete Alignment einen
vorgegebenen Schwellwert überschreitet -- nützlich für die
Leistungsanalyse realer Betriebssystemkerne.

\subsubsection{Parallele und verteilte Berechnung des Subgraph Algorithmus}

Die $n$ Rotationsschritte des Subgraph Algorithmus sind gemäß
Lemma~\ref{lem:unabhaengigkeit} voneinander unabhängig.
Dies ermöglicht eine \emph{embarrassingly parallel} Implementierung auf
Mehr\-kern\-pro\-zes\-so\-ren: Jeder der $n$ Rotationsschritte wird einem
separaten CPU-Kern zugewiesen; die $n$ LCS-Berechnungen laufen simultan.
Die Gesamtlaufzeit sinkt von $O(n^3)$ auf $O(n^2)$ bei $n$ verfügbaren
Kernen~\cite{openmp2011}.
Darüber hinaus lässt sich die DP-Matrix jeder LCS-Berechnung auf SIMD-fähige
Vektorregister (AVX-512) abbilden, was die Berechnungslatenz pro Rotation
um einen weiteren konstanten Faktor reduziert~\cite{rognes2011faster}.
Für extrem große Graphen ($n > 10^4$) ist eine GPU-basierte Implementierung
(CUDA/OpenCL) denkbar, bei der jede Zeile der DP-Tabelle als separater
CUDA-Thread verarbeitet wird.

\subsubsection{Anwendung auf formale Verifikation von Synchronisationsprotokollen}

In der formalen Methodik werden Synchronisationsprotokolle häufig als
Kripke-Strukturen oder Transitionssysteme modelliert und mittels
Model-Checking auf Eigenschaften wie Deadlock-Freiheit und
starvation-Freiheit überprüft~\cite{clarke1999model}.
Der Subgraph Algorithmus kann hier als \emph{Vorstufe} eingesetzt werden:
Bevor ein teures Model-Checking durchgeführt wird, prüft er, ob zwei
Synchronisationsprotokolle strukturell verwandt sind. Falls
$G_A \subseteq G_B$ gilt und $G_B$ bereits als deadlock-frei verifiziert
wurde, lässt sich diese Eigenschaft unter zusätzlichen Monotonie-Bedingungen
auf $G_A$ übertragen -- eine Technik, die als
\emph{compositional reasoning}~\cite{mcmillan1993symbolic} bekannt ist.

\subsubsection{Integration in statische Analysetools und Compiler}

Moderne Compiler (GCC, LLVM/Clang) erzeugen interne Darstellungen von
nebenläufigem Code, die als Kontrollflussgraphen (CFG) oder
Threaded-Code-Graphen vorliegen~\cite{lattner2004llvm}. Der Subgraph
Algorithmus könnte als Compiler-Pass implementiert werden, der erkennt,
wenn ein Programmierer einen Mutex-Lock verwendet, obwohl ein einfacher
Spin-Lock ausreichen würde -- oder umgekehrt, wenn ein Spin-Lock
einen Mutex ersetzen sollte (z.\,B.\ bei langen kritischen Regionen).
Diese automatische Empfehlung würde die Systemleistung verbessern, ohne
dass Programmierer selbst fundierte Kenntnisse aller Synchronisations\-mechanismen
benötigen.

\subsubsection{Erweiterung auf hierarchische Parallelität (OpenMP, MPI)}

Die hier betrachteten Synchronisationsmechanismen beziehen sich
hauptsächlich auf POSIX-Threads (Pthreads) innerhalb eines Prozesses.
In modernen Hochleistungsrechnersystemen spielen jedoch auch
Multi-Ebenen-Parallelitätsmodelle wie OpenMP~\cite{openmp2011} (für
Shared-Memory-Parallelismus) und MPI~\cite{mpi2009} (für
Distributed-Memory-Parallelismus) eine wesentliche Rolle. Die
Zustandsgraphen von MPI-Kollektiv-Operationen (z.\,B.\
\texttt{MPI\_Barrier}, \texttt{MPI\_Reduce}) könnten analog zu den
hier definierten Synchronisationsgraphen modelliert und mittels
Subgraph Algorithmus klassifiziert werden. Eine Hierarchie
$G_{\mathrm{Barrier}} \subseteq G_{\mathrm{AllReduce}}$ würde
beispielsweise belegen, dass eine Barriere strukturell in einer
AllReduce-Operation enthalten ist -- mit direkten Implikationen für
die Optimierung kollektiver MPI-Kommunikation.

\subsubsection{Maschinelles Lernen zur Subgraph-Vorhersage}

Graph Neural Networks (GNNs)~\cite{scarselli2009graph} wurden in den
letzten Jahren erfolgreich für Subgraph-Erkennungsaufgaben eingesetzt.
Eine interessante Forschungsrichtung besteht darin, den Subgraph Algorithmus
als \emph{Label} für das Training eines GNN zu verwenden: Für eine große
Datenbank von Synchronisationsgraphen werden die Subgraph-Beziehungen
mittels Subgraph Algorithmus vorberechnet; das GNN lernt anschließend, diese
Beziehungen in $O(1)$ vorherzusagen (nach dem Training).
Dieser Ansatz kombiniert die Korrektheit des exakten Algorithmus mit der
Schnelligkeit eines gelernten Modells -- besonders interessant für
Echtzeit-Anwendungen wie die Live-Analyse laufender Betriebssystem\-prozesse.

\subsubsection{Erweiterung auf hybride Synchronisationsprotokolle}

In der Praxis verwenden Betriebssysteme hybride Protokolle, die mehrere
Mechanismen kombinieren: Der Linux-Kernel nutzt z.\,B.\ \texttt{futex}
(Fast User-space Mutex), der im Unstrittenen Fall im Userspace operiert
und nur bei Konkurrenz in den Kernel wechselt~\cite{franke2002fuss}.
Der Zustandsgraph eines Futex besitzt daher Teile des Spin-Lock-Graphen
(für den streitfreien Fall) und Teile des Mutex-Lock-Graphen (für den
strittigen Fall). Eine Erweiterung des Subgraph Algorithmus auf
\emph{zusammengesetzte} Graphen, die als Vereinigung von Teilgraphen
verschiedener Mechanismen modelliert werden, wäre ein wertvoller Beitrag.
Dies erfordert eine Anpassung der Signaturberechnung auf Teilgraph-Mengen
und eine modifizierte LCS-Variante, die Teilsequenzen aus mehreren
Quellgraphen berücksichtigt.

\subsubsection{Anwendung auf Echtzeitsysteme und AUTOSAR}

In eingebetteten Echtzeitsystemen (RTOS) und dem AUTOSAR-Standard~\cite{autosar2021}
spielen Synchronisationsmechanismen eine sicherheitskritische Rolle.
Prioritätsinversion -- das bekannte Problem, bei dem ein niedrigpriorer
Prozess einen hochprioren blockiert -- ist direkt als Subgraph-Eigenschaft
formulierbar: Tritt der Subgraph $G_{\mathrm{priorityInversion}} \subseteq
G_{\mathrm{Protokoll}}$ auf, ist das Protokoll anfällig für dieses Problem.
Die Anwendung des Subgraph Algorithmus auf AUTOSAR-Kommunikations\-protokolle
könnte automatisch solche sicherheitskritischen Muster erkennen, bevor ein
Steuergerät in die Serienfertigung geht~\cite{iso26262}.

\subsubsection{Quantencomputing und Quanten-Synchronisation}

Mit dem Aufkommen von Quantencomputern entstehen neue Parallelitätsparadigmen,
in denen klassische Konzepte wie Mutex-Locks und Semaphoren nicht direkt
übertragbar sind~\cite{gay2006quantum}. Dennoch existieren Ansätze für
Quanten-Zustandsautomaten, die ähnliche Strukturen wie die hier modellierten
Synchronisationsgraphen aufweisen. Eine Forschungsfrage ist, ob der Subgraph
Algorithmus auf Quanten-Zustandsgraphen -- modelliert als unitäre
Transformationen auf Hilbert-Räumen -- erweitert werden kann. Dies würde
erfordern, die diskrete Signatur-Funktion durch eine quantenmechanische
Observable zu ersetzen und die LCS-Berechnung auf Quanten-Parallelismus
zu verallgemeinern.

\subsubsection{Offene Komplexitätsfragen}

Aus theoretischer Sicht bleibt die Frage offen, ob eine \emph{unbedingte}
untere Schranke von $\Omega(n^3)$ ohne die SETH-Hypothese beweisbar ist.
Aktuelle Ergebnisse~\cite{abboud2015tight} zeigen lediglich eine bedingte
Schranke unter SETH. Eine Auflösung dieser Frage wäre ein bedeutender
Fortschritt in der Komplexitätstheorie des Subgraph-Isomorphismus.
Darüber hinaus ist unklar, ob der Subgraph Algorithmus auf
\emph{approximative} Subgraph-Beziehungen (LCS $\geq \alpha \cdot n$ für
ein festes $\alpha \in (0, 1)$) in besserer Laufzeit als $O(n^3)$
verallgemeinert werden kann, ohne SETH zu widerlegen.

% =========================================================================
\begin{thebibliography}{99}
% =========================================================================

\bibitem{epp2026subgraph}
  S.~Epp,
  \textit{Der Subgraph Algorithmus},
  Technischer Bericht, Bielefeld, Januar 2026.
  \url{https://gitlab.com/epp-group/subgraph}

\bibitem{tanenbaum2014modern}
  A.~S.~Tanenbaum und H.~Bos,
  \textit{Modern Operating Systems}, 4.~Auflage,
  Pearson, 2014.

\bibitem{silberschatz2018os}
  A.~Silberschatz, P.~B.~Galvin und G.~Gagne,
  \textit{Operating System Concepts}, 10.~Auflage,
  Wiley, 2018.

\bibitem{herlihy2008art}
  M.~Herlihy und N.~Shavit,
  \textit{The Art of Multiprocessor Programming},
  Morgan Kaufmann, 2008.

\bibitem{love2010linux}
  R.~Love,
  \textit{Linux Kernel Development}, 3.~Auflage,
  Addison-Wesley, 2010.

\bibitem{smith1981identification}
  T.~F.~Smith und M.~S.~Waterman,
  ``Identification of Common Molecular Subsequences'',
  \textit{Journal of Molecular Biology},
  Bd.~147, S.~195--197, 1981.

\bibitem{openmp2011}
  OpenMP Architecture Review Board,
  \textit{OpenMP Application Program Interface Version 3.1},
  2011.
  \url{https://www.openmp.org}

\bibitem{rognes2011faster}
  T.~Rognes,
  \enquote{Faster Smith-Waterman Database Searches with Inter-Sequence
  SIMD Parallelisation},
  \textit{BMC Bioinformatics}, 12:221, 2011.

\bibitem{clarke1999model}
  E.~M.~Clarke, O.~Grumberg und D.~Peled,
  \textit{Model Checking},
  MIT Press, 1999.

\bibitem{mcmillan1993symbolic}
  K.~L.~McMillan,
  \textit{Symbolic Model Checking},
  Kluwer Academic Publishers, 1993.

\bibitem{lattner2004llvm}
  C.~Lattner und V.~Adve,
  \enquote{LLVM: A Compilation Framework for Lifelong Program Analysis
  and Transformation},
  \textit{Proc.\ CGO~2004}, S.~75--86.

\bibitem{mpi2009}
  MPI~Forum,
  \textit{MPI: A Message-Passing Interface Standard, Version 2.2},
  2009.
  \url{https://www.mpi-forum.org}

\bibitem{scarselli2009graph}
  F.~Scarselli et al.,
  ``The Graph Neural Network Model'',
  \textit{IEEE Transactions on Neural Networks},
  Bd.~20, Nr.~1, S.~61--80, 2009.

\bibitem{franke2002fuss}
  H.~Franke, R.~Russell und M.~Kirkwood,
  ``Fuss, Futexes and Furwocks: Fast Userland Locking in Linux'',
  \textit{Proc.\ Ottawa Linux Symposium}, S.~479--495, 2002.

\bibitem{autosar2021}
  AUTOSAR Consortium,
  \textit{AUTOSAR Classic Platform -- Release 4.4.0},
  2021.
  \url{https://www.autosar.org}

\bibitem{iso26262}
  ISO,
  \textit{ISO 26262: Road Vehicles -- Functional Safety},
  2nd~edition, 2018.

\bibitem{gay2006quantum}
  S.~J.~Gay,
  ``Quantum Programming Languages: Survey and Bibliography'',
  \textit{Mathematical Structures in Computer Science},
  Bd.~16, Nr.~4, S.~581--600, 2006.

\bibitem{abboud2015tight}
  A.~Abboud, A.~Backurs und V.~V.~Williams,
  \enquote{Tight Hardness Results for LCS and Other Sequence Similarity
  Measures},
  \textit{Proc.\ FOCS~2015}, S.~59--78.

\end{thebibliography}

\end{document}
